% !TEX root = ../notes.tex

\documentclass[../notes.tex]{subfiles}

\begin{document}

\section{May 6}
Today, we define $L$-functions.

\subsection{Multiplicity One, Finally}
We are now ready to sketch the proof of multiplicity one.
\begin{theorem}[Multiplicity one]
	Fix a cuspidal automorphic form $\pi$ for the group $\op{GL}_n$ over a function field $F$. Then $\pi$ is uniquely embedded in $\mc A_{\mathrm{cusp}}$.
\end{theorem}
We will only discuss the proof for $n\in\{2,3\}$.
\begin{proof}[Proof for $n=2$]
	There are two steps: we will show that $\pi$ has a Whittaker model, and then we will show that the Whittaker model actually determines $\pi$.

	Let's start by hunting for a Whittaker functional. Roughly speaking, we see that we are trying to define an embedding $\mc A_{\mathrm{cusp}}\to\mc A((N(\AA_F),\psi_N)\backslash G(\AA_F))$. We start with some $\varphi\in V_\pi$, which is lives in the space of functions $G(F)\backslash G(\AA_F)\to\CC$, and we are hunting for an invariant function for $N(\AA_F)$. Because $\psi_N$ already vanishes on $N(F)$, we are thus allowed to do the average
	\[W_\bullet(g)\coloneqq\int_{N(F)\backslash N(\AA_F)}\varphi(ug)\psi(u)^{-1}\,du.\]
	Notably, the space $N(F)\backslash N(\AA_F)$ is an iterated extension of $F\backslash\AA_F$, so it is compact, so this integral converges. Now, to show that we actually land in $\mc A((N(\AA_F),\psi_N)\backslash G(\AA_F))$, there are some extra checks on smoothness and so on which are required, but we will skip them.
	% One needs to check that $\Lambda$ is nonzero, which uses cuspidality. From this functional $\Lambda$, there is a map $V_\pi\to\mc A((N(\AA_F),\psi_N)\backslash G(\AA_F))$ given by sending $\varphi$ to the function $g\mapsto\Lambda(\pi(g)\varphi)$. Notably, there are some extra checks one has to do in the number field case, such as moderate growth, which we will skip.

	Let's next discuss why $W_\bullet$ is actually injective. The idea is to take a Fourier transform. Indeed, note that we can replace $\psi$ in the definition of $W_\bullet$ with any other additive character. Recall that the Pontryagin dual of $F\backslash\AA_F$ is $F$, so other additive characters look like $\psi_a(x)\coloneqq\psi(ax)$. Accordingly, we define
	\[W_\varphi(g;\psi_a)\coloneqq\int_{N(F)\backslash N(\AA_F)}\varphi(ug)\psi_a(u)^{-1}\,du.\]
	For example, for $G=\op{GL}_2$, we find that
	\[W_\varphi(g;\psi_a)=\int_{N(F)\backslash N(\AA_F)}\varphi\left(\begin{bmatrix}
		1 & a^{-1}x \\ & 1
	\end{bmatrix}g\right)\psi(-x)\,dx.\]
	Using the invariance of $\varphi$, one can write $\begin{bsmallmatrix}
		1 & a^{-1}x \\ & 1
	\end{bsmallmatrix}=\op{diag}(1/a,1)\begin{bsmallmatrix}
		1 & x \\ & 1
	\end{bsmallmatrix}\op{diag}(a,1)$, so we see that
	\begin{equation}
		W_\varphi(g;\psi_a)=W_\varphi(\op{diag}(a,1)g). \label{eq:relate-whittakers}
	\end{equation}
	Now, we apply Fourier inversion. View $f_{\varphi,g}\colon x\mapsto\varphi\left(\begin{bsmallmatrix}
		1 & x \\ & 1
	\end{bsmallmatrix}g\right)$ as a function on $F\backslash\AA_F$. Notably, the Fourier transform $\widehat f_{\varphi,g}(\psi')$ is $W_\varphi(g;\psi')$ by construction. (Here, the Fourier transform takes as input an element $\psi'$ of the Pontryagin dual $(F\backslash\AA_F)^*$.) Thus, Fourier inversion implies
	\[f_{\varphi,g}(0)=\sum_{\psi'\in(F\backslash\AA_F)^*}\widehat f_{\varphi,g}(\psi').\]
	Note that the left-hand side is $\varphi(g)$, so the moral is that
	\[\varphi(g)=\sum_{\psi'\in(F\backslash\AA_F)^*}W_\varphi(g;\psi').\]
	In particular, we can recover $\varphi$ from its image in the Whittaker functional! Now, using our parameterization of $(F\backslash\AA_F)^*$, we see that
	\[\varphi(g)=W_\varphi(g;1)+\sum_{a\in F^\times}W_\varphi(g;\psi_a),\]
	which is
	\[\varphi(g)=W_\varphi(g;1)+\sum_{a\in F^\times}W_\varphi(\op{diag}(a,1)g;\psi)\]
	by \eqref{eq:relate-whittakers}. We are now ready to show that $W_\bullet$ is injective: if $W_\varphi=0$, then $\varphi(g)=W_\varphi(g;1)$, but this is
	\[W_\varphi(g;1)=\int_{F\backslash\AA_F}\varphi\left(\begin{bmatrix}
		1 & x \\ & 1
	\end{bmatrix}\right)\,dg.\]
	This vanishes because $\varphi$ is cuspidal! In fact, we can see that if $\varphi$ is cuspidal, then this argument shows that
	\[\varphi(g)=\sum_{a\in F^\times}W_\varphi(\op{diag}(a,1)g).\]
	This establishes that any cuspidal automorphic representation $\pi$ admits a nonzero Whittaker functional.

	It remains to recover $\pi$ from its Whittaker model. Note that $\op{Hom}_{G(\AA_F)}(\pi,\mc A_{\mathrm{cusp}})$ embeds into
	\[\op{Hom}_{G(\AA_F)}(\pi,\mc A((N(\AA_F),\psi_N)\backslash G(\AA_F))).\]
	By Frobenius reciprocity, this embeds into $\op{Hom}_{N(\AA_F)}(\pi,\psi_N)$, and as discussed in \Cref{thm:global-whittaker-unique}, it is an isomorphism onto the subspace of moderate growth. However, the subspace of moderate growth is simply the space of Whittaker functionals, so we are done.
	% We will show that $\op{Hom}(\pi,\mc A_{\mathrm{cusp}})$ embeds into the space of Whittaker functionals of $\pi$; this will complete the proof by \Cref{thm:global-whittaker-unique}. Well, for any map $i\colon\pi\to\mc A_{\mathrm{cusp}}$, and consider the composite $\Lambda_i$ defined by
	% \[\pi\stackrel i\to\mc A_{\mathrm{cusp}}\stackrel W\to\mc A((N(\AA_F),\psi_N)\backslash G(\AA_F))\stackrel{\op{ev}_1}\to\CC.\]
	% Note that $\Lambda_i$ is a Whittaker functional by construction, so we would like to show that the map $i\mapsto\Lambda_i$ is injective. Well, suppose that $\Lambda_i=0$, which means that $W_{i(v)}(1)=0$ for all $v\in\pi$. We would like to show that $i(v)=0$; for brevity, set $\varphi\coloneqq i(v)$. For any $g\in G(\AA_F)$, we see that $R_g\varphi$ continues to be in $\im i$, so $W_{R_g\varphi}(1)=0$, so $W_\varphi(g)=0$ for all $g$, so $\varphi=0$.
\end{proof}
\begin{proof}[Sketch for $n=3$]
	Here, we cannot directly take a Fourier transform over $N(\AA_F)$, which is not abelian. However, we could take a Fourier transform over the abelian subgroup
	\[U_2\coloneqq\left\{\begin{bmatrix}
		1 & & a \\ & 1 & b \\ && 1
	\end{bmatrix}:a,b\in\mathbb G_a\right\}.\]
	Now, given $\varphi\in\mc A$, we define $\varphi_1(g)\coloneqq\int_{U_2(F)\backslash U_2(\AA_F)}\varphi(ug)\psi(-u)\,du$. Notably, $\varphi_1$ is the Fourier transform of the function
	\[(y,x)\mapsto\varphi\left(\begin{bmatrix}
		1 && y \\ & 1 & x \\ && 1
	\end{bmatrix}g\right)\]
	evaluated at the character $(0,\psi)\in((F\backslash\AA_F)^*)^2$. Now, calculating as before, one can recover the value of the Fourier transform at $(\psi_1,\psi_2)\ne(0,0)$ from the values of $\varphi_1\left(\op{diag}(\gamma,1)g\right)$ as $\gamma$ varies over $\op{GL}_2(F)$. Additionally, the same sort of calculation from $n=2$ finds that $\varphi$ can be recovered from $\varphi_1$ by Fourier inversion, provided that $\varphi$ is cuspidal!

	Technically, the previous paragraph only uses that $\varphi$ is invariant under the subgroup $Q_3(F)$, where $Q_3=\mathrm{GL}_2U_2$, and one can see that $\varphi_1$ is then invariant under $Q_2(F)$, where $Q_2(F)$ is generated by $\op{diag}(a,1,1)$ and $\op{diag}\left(\begin{bsmallmatrix}
		1 & x \\ & 1
	\end{bsmallmatrix},1\right)$. However, there is some residual cuspidality to use on $\varphi_1$, and applying the same sort of argument to the construction
	\[\varphi_2(g)\coloneqq\int_{F\backslash\AA_F}\varphi_1\left(\begin{bmatrix}
		1 & x \\ & 1 \\ && 1
	\end{bmatrix}g\right)\psi(-x)\,dx.\]
	Redoing the argument, we find that $\varphi_2$ lives in $\mc A((N(\AA_F),\psi_N)\backslash\op{GL}_3(\AA_F))$, and one can make the same sort of Whittaker argument as before.
\end{proof}

\subsection{Global Zeta Integrals}
We are now ready to define some $L$-functions. For motivation, recall that Tate's thesis defined a global zeta integral
\[Z(f,\chi)\coloneqq\int_{\AA_F^\times}f(a)\chi(a)\,d^\times a,\]
where $f$ is some Schwartz function on $\AA_F$, and $\chi$ is a continuous character on $F^\times\backslash\AA_F^\times$. This integral had the following notably properties.
\begin{itemize}
	\item If $f$ is a pure tensor, then this unfolds into a product of local zeta integrals.
	\item For given $f$ and $\chi$, the function $Z(s,f,\chi)\coloneqq Z\left(f,\chi\left|\cdot\right|^s\right)$ has a meromorphic continuation with well-understood poles.
	\item There was a functional equation $Z(f,\chi)=Z\left(\hat f,\chi^{-1}\left|\cdot\right|\right)$.
\end{itemize}
We would like to build an integral for $\op{GL}_2$ with similar properties.
\begin{definition}[global integral]
	Fix a global field $F$, and choose a cuspidal automorphic representation $\pi$ of $\op{GL}_2(\AA_F)$. Given a character $\chi\colon F^\times\backslash\AA_F^\times\to\CC^\times$, we define
	\[Z(\varphi,\chi)\coloneqq\int_{F^\times\backslash\AA_F^\times}\varphi\left(\begin{bmatrix}
		a \\ & 1
	\end{bmatrix}\right)\chi(a)\,d^\times a.\]
\end{definition}
\begin{remark}
	Here, the ``main input'' is $\varphi$, and $\chi$ is an auxiliary ``test function.'' Namely, $\chi$ is akin to $f$ in Tate's original integral. The correct choice of test function is predicted by the Braverman--Kazhdan program.
\end{remark}
\begin{remark}
	Perhaps one should check that this integral converges. For fixed $\chi$ and $\varphi$, one can define
	\[Z(s,\varphi,\chi)\coloneqq\int_{F^\times\backslash\AA_F^\times}\varphi\left(\begin{bmatrix}
		a \\ & 1
	\end{bmatrix}\right)\chi(a)\left|a\right|^{s-1/2}\,d^\times a.\]
	It turns out that this is a holomorphic function in $s\in\CC$, which we will show shortly.
\end{remark}
To be able to say anything interesting about our zeta integral, we should use the Whittaker model $W_\bullet$ of $\pi$. For example, let's explain how we get an Euler product.
\begin{lemma}
	Fix a global field $F$, and choose a cuspidal automorphic representation $\pi$ of $\op{GL}_2(\AA_F)$. If $\varphi\in\pi$ is a pure tensor, then $Z(\varphi,\chi)$ admits an Euler product expansion.
\end{lemma}
\begin{proof}
	The problem is that we are integrating over $F^\times\backslash\AA_F^\times$, which does not split as a product. Thus, we need to ``sum over $F^\times$'' to be able to divide up our integral. Namely, for $\varphi\in\pi$, one has
	\[\varphi(g)=\sum_{b\in F^\times}W_\varphi\left(\begin{bmatrix}
		b \\ & 1
	\end{bmatrix}g\right).\]
	In particular,
	\[\varphi(\op{diag}(a,1))=\sum_{b\in F^\times a}W_\varphi(\op{diag}(b,1)),\]
	so
	\[Z(\varphi,\chi)=\int_{F^\times\backslash\AA_F^\times}\Bigg(\sum_{b\in F^\times a}W_\varphi(\op{diag}(b,1))\Bigg)\chi(a)\,d^\times a.\]
	Now, one can combine the sum and integral into
	\[Z(\varphi,\chi)=\int_{\AA_F^\times}W_\varphi(\op{diag}(b,1))\chi(b)\,d^\times b,\]
	which we can now factor.
	\begin{remark}
		This step is known as ``unfolding.''
	\end{remark}
	To explain the factorization, write $\pi=\bigotimes_v'\pi_v$, and each $\pi_v$ admits a Whittaker model $W_v$, and we choose scalars appropriately so that $W=\bigotimes_vW_v$. Now if $\varphi$ is a pure tensor $\bigotimes\varphi_v$. Thus, if $\varphi$ is a pure tensor $\bigotimes_v\varphi_v$, we see that
	\[Z(\varphi,\chi)=\prod_v\left(\int_{F^\times}W_{\varphi_v}(b_v)\chi_v(b_v)\,d^\times b_v\right).\]
	Notably, viewing $W_{\varphi_v}$ as some one-variable function, this is basically one of Tate's local integral.
\end{proof}
\begin{remark}
	One can show convergence of the integral from the above discussion. It turns out that the function $a\mapsto\varphi\left(\op{diag}(a,1)g\right)$ is rapidly decreasing (namely, faster than inverses of polynomials) as $\left|a\right|\to\infty$. (Indeed, we have seen this already for cuspidal modular forms.) Using the Weyl element $\begin{bsmallmatrix}
		& 1 \\ 1
	\end{bsmallmatrix}$ (and fixing a central character for $\pi$), we see that the same is true as $\left|a\right|\to0$. Technically, one needs to know that the central character is unitary, which we can do by twisting $\pi$ by the determinant. Convergence of the integral follows by splitting the global zeta integral into the regimes $\left|a\right|\ge1$ and $\left|a\right|\le1$.
\end{remark}
\begin{remark}
	Different choices of Whittaker models may change individual terms in the Euler product expansion.
\end{remark}
Let's now explain how one could do some calculations. Fix a place $v$ where $\pi_v$ is spherical, and let $\varphi_v$ be a nonzero spherical vector. Then we would like to calculate
\[\int_{F_v^\times}W_{\varphi_v}(\op{diag}(a,1))\left|a\right|^{s-1/2}\,d^\times a.\]
One can explicitly calculate $W_{\varphi_v}$ under the assumption that $\varphi_v$ is spherical, which allows one to compute this.

\end{document}